\documentclass[12pt, a4paper]{article}
\usepackage[utf8]{inputenc}
\usepackage{geometry}
\usepackage{hyperref}
\usepackage{graphicx}
\usepackage{listings}
\usepackage{xcolor}

\geometry{margin=1in}

\title{Pharma-Cure V4: Unified Clinical Network \\ \large Technical Specification \& Architecture}
\author{Pharma-Cure Engineering Team}
\date{June 2026}

\begin{document}

\maketitle

\section{Introduction}
Pharma-Cure V4 is an advanced, geolocation-aware pharmacy management ecosystem. Unlike traditional inventory systems, V4 implements a dynamic \textbf{Proposal-based Economy} (Uber-like model) where patients, pharmacists, and suppliers negotiate prices and quantities in real-time.

\section{System Architecture}
The system is built on a Node.js Express backend with a MySQL relational database. The frontend is a sophisticated Single Page Application (SPA) utilizing CSS Grid and modern JavaScript.

\section{Economic Model: The Proposal System}
Fixed price lists have been deprecated in favor of a \textbf{Negotiated Proposal} model.
\begin{itemize}
    \item \textbf{Patient to Pharmacy:} Patients broadcast a medicine requirement with a target price. Pharmacies accept or reject based on stock and profitability.
    \item \textbf{Pharmacy to Supplier:} Pharmacists propose restock batches to suppliers. The system provides a ``System Rating'' (Excellent, Good, Fair, Poor) to advise on market competitiveness.
\end{itemize}

\section{Database Schema}
The relational model consists of five primary entities: Users, Patients, Pharmacies, Medicines, and Proposals. Geolocation data (latitude and longitude) is stored with 8-decimal precision to ensure high accuracy.

\section{UI/UX Philosophy}
The visual identity follows a \textbf{Clean Clinical} aesthetic. The design is optimized for high-density information display, removing non-professional elements such as emojis and utilizing the full screen width (wrap-around layout) for complex data tables and analytics.

\end{document}
